\documentclass[11pt]{article}
\usepackage[margin=1in]{geometry}
\usepackage{amsmath,amssymb,amsthm}
\usepackage{booktabs}
\usepackage[hidelinks]{hyperref}

\newtheorem{proposition}{Proposition}
\newtheorem{corollary}{Corollary}[proposition]
\theoremstyle{definition}
\newtheorem{definition}{Definition}
\theoremstyle{remark}
\newtheorem*{remark}{Remark}

\newcommand{\A}{A(a)}
\newcommand{\Sib}{S(a)}
\newcommand{\w}{w}
\newcommand{\anc}{\mathrm{anc}}
\newcommand{\full}{\mathrm{full}}
\newcommand{\win}{\mathrm{win}}

\title{Ancestor-Only Context Selection: Formal Framework}
\author{Okechukwu Zeal Achonu}
\date{}

\begin{document}
\maketitle

\section{Setting}

\begin{definition}[Branched session]
A session is a rooted tree $T=(V,E)$ of messages with root $r$. Every
$v\neq r$ has a parent $p(v)$. A logical clock $\tau:V\to\mathbb{N}$ is
injective and increases along every edge: $\tau(p(v))<\tau(v)$. Each message
has a token weight $\w(v)>0$ under a fixed reference tokenizer, and for
$U\subseteq V$ we write $\w(U)=\sum_{u\in U}\w(u)$.
\end{definition}

\begin{definition}[Active path and siblings]
The user asks a question $q$ as a child of an \emph{ask node} $a\in V$, at time
$\tau_q>\max_{v\in V}\tau(v)$. The \emph{active path} is the set of $a$ and
its ancestors,
\[
  \A=\{a,\,p(a),\,p(p(a)),\,\dots,\,r\},
\]
and the \emph{sibling content} is $\Sib=V\setminus\A$.
\end{definition}

\begin{definition}[Context rules]
Each rule returns a set of messages, sent in the stated order, followed by $q$.
\begin{itemize}
  \item Ancestor-only: $C_\anc=\A$, ordered root to node.
  \item Full tree: $C_\full=\{v\in V:\tau(v)<\tau_q\}=V$, ordered by $\tau$.
  \item Sliding window with budget $B=\w(\A)$: list $V$ by decreasing $\tau$ as
        $v_1,v_2,\dots,v_{|V|}$ and let $C_\win=\{v_1,\dots,v_m\}$ with $m$ the
        largest index such that $\sum_{i\le m}\w(v_i)\le B$.
\end{itemize}
The window is a \emph{suffix} of the session in time, it never splits a
message, and it is \emph{maximal}: if $m<|V|$ then $\w(C_\win)+\w(v_{m+1})>B$.
\end{definition}

\section{Token efficiency}

\begin{proposition}[Containment and efficiency]\label{prop:eff}
$C_\anc\subseteq C_\full$, so $\w(C_\anc)\le\w(C_\full)$, and the token
efficiency of ancestor-only over full tree equals the siblings' share of the
tree:
\[
  E \;=\; 1-\frac{\w(C_\anc)}{\w(C_\full)} \;=\; \frac{\w(\Sib)}{\w(V)} .
\]
\end{proposition}
\begin{proof}
Every element of $\A$ is in $V$, so $C_\anc=\A\subseteq V=C_\full$, and
$\w$ is a sum of positive terms, hence monotone. Since $V=\A\,\dot\cup\,\Sib$
is a disjoint union, $\w(V)=\w(\A)+\w(\Sib)$, so
$1-\w(\A)/\w(V)=\w(\Sib)/\w(V)$.
\end{proof}

\begin{remark}
$E$ needs no model; it follows from the tree shape alone. In the frozen
dataset the measured $E$ equals $\w(\Sib)/\w(V)$ to machine precision in
every scenario, as the proposition requires.
\end{remark}

\begin{proposition}[Growth over a session]\label{prop:growth}
Suppose a session grows one message at a time, $V_1\subset V_2\subset\cdots\subset V_n$
with $|V_t|=t$, and each new message triggers one model call whose prompt is
the current context. Let $0<w_{\min}\le\w(v)\le w_{\max}$ and let $d_t$ be the
depth of the node being answered at step $t$. Then the cumulative prompt cost
satisfies
\[
  \sum_{t=1}^{n}\w(V_t)\;\ge\;w_{\min}\,\frac{n(n+1)}{2}
  \qquad\text{and}\qquad
  \sum_{t=1}^{n}\w\bigl(A_t\bigr)\;\le\;w_{\max}\sum_{t=1}^{n}(d_t+1).
\]
If depth is bounded, $d_t\le D$, full-tree cost grows as $\Theta(n^2)$ while
ancestor-only cost grows as $O(nD)$, i.e.\ linearly in $n$.
\end{proposition}
\begin{proof}
$\w(V_t)\ge t\,w_{\min}$ because $|V_t|=t$; summing gives the first bound.
$|A_t|=d_t+1$ because a root-to-node path at depth $d_t$ has $d_t+1$ nodes,
so $\w(A_t)\le (d_t+1)w_{\max}$; summing gives the second. With $d_t\le D$
the right-hand side is at most $nw_{\max}(D+1)$. The full-tree sum is also at
most $w_{\max}n(n+1)/2$, so it is $\Theta(n^2)$.
\end{proof}

\section{Leakage is only possible through exposure}

\begin{definition}[Sibling-only fact]
A fact $f$ (a string, e.g.\ a canary) is \emph{sibling-only} if every message
whose text contains $f$ lies in $\Sib$, and $f$ does not occur in $q$.
\end{definition}

\begin{proposition}[Ancestor-only cannot contain sibling facts]\label{prop:nosib}
If $f$ is sibling-only, then $f$ does not occur in the prompt of $C_\anc$.
\end{proposition}
\begin{proof}
The prompt of $C_\anc$ is the text of the messages in $\A$ followed by $q$.
No message of $\A$ contains $f$ (all such messages lie in $\Sib$, which is
disjoint from $\A$), and $q$ does not contain $f$.
\end{proof}

\begin{definition}[Exposure, chance rate, uptake]
For a rule $c$, let $X_c(f)=\mathbf 1[f\text{ occurs in the prompt of }C_c]$.
Let $M$ be the event that the model's answer contains $f$. Let
$\pi_f=\Pr(M\mid f\text{ absent from the prompt})$ be the chance rate, and
$u_c(f)=\Pr(M\mid X_c(f)=1)$ the uptake under rule $c$.
\end{definition}

\begin{proposition}[Exposure--uptake decomposition]\label{prop:decomp}
$X_c(f)$ is a deterministic function of $(T,\tau,\w,a)$ and can be computed
without calling a model. For a sibling-only $f$, the excess mention rate of
rule $c$ over ancestor-only is
\[
  \Pr(M\mid c)-\Pr(M\mid\anc)\;=\;X_c(f)\,\bigl(u_c(f)-\pi_f\bigr).
\]
\end{proposition}
\begin{proof}
$C_c$ is defined from $T$, $\tau$, $\w$ and $a$ alone (Definition 3), so
whether $f$ occurs in its text is determined by them. By
Proposition~\ref{prop:nosib}, $f$ is absent under ancestor-only, so
$\Pr(M\mid\anc)=\pi_f$. Under rule $c$, if $X_c(f)=0$ the prompt does not contain
$f$ and $\Pr(M\mid c)=\pi_f$; if $X_c(f)=1$ then $\Pr(M\mid c)=u_c(f)$. In both
cases $\Pr(M\mid c)-\pi_f=X_c(f)(u_c(f)-\pi_f)$.
\end{proof}

\begin{remark}
Two consequences shape the measurement. (i) For random canaries
$\pi_f\approx 0$: in the full run, 0 of 2,598 ancestor-only answers contained a
canary. For forced-choice values (e.g.\ ``LRU''), $\pi_f$ can be large and must
be measured on ancestor-only answers; 17 of 2,160 were chance hits. (ii) The
experiment therefore logs $X_c(f)$ for every call and estimates leakage among
exposed calls, which is the uptake $u_c(f)$.
\end{remark}

\section{What the sliding window displaces}

\begin{proposition}[Membership test]\label{prop:member}
A message $m$ is in $C_\win$ if and only if
\[
  \sum_{u\in V:\ \tau(u)\ge\tau(m)} \w(u)\;\le\;B .
\]
\end{proposition}
\begin{proof}
$C_\win$ is the longest prefix of the list ordered by decreasing $\tau$ whose
total weight is at most $B$. Since all weights are positive, the prefix sums
increase strictly, so $m$ belongs to that prefix exactly when the prefix ending
at $m$, which is $\{u:\tau(u)\ge\tau(m)\}$, has weight at most $B$.
\end{proof}

\begin{corollary}[The root never survives]\label{cor:root}
If $\Sib\neq\varnothing$, the root $r$ is not in $C_\win$.
\end{corollary}
\begin{proof}
$\tau(r)$ is the minimum of $\tau$ over $V$, so the sum in
Proposition~\ref{prop:member} for $m=r$ is
$\w(V)=\w(\A)+\w(\Sib)=B+\w(\Sib)>B$.
\end{proof}

\begin{proposition}[One-for-one displacement]\label{prop:displace}
Let $\sigma=\w(C_\win\cap\Sib)$ be the sibling tokens admitted to the window.
Then
\[
  \w(\A)-\w(C_\win\cap\A)\;\ge\;\sigma ,
\]
so every admitted sibling token displaces at least one active-path token.
Moreover $C_\win\cap\A$ is a suffix of the path in time: the active messages
lost are always the oldest, root-side ones.
\end{proposition}
\begin{proof}
$\w(C_\win\cap\A)+\sigma=\w(C_\win)\le B=\w(\A)$; rearrange. For the second
claim, if $m\in\A$ is in the window and $m'\in\A$ has $\tau(m')>\tau(m)$, then
the sum in Proposition~\ref{prop:member} for $m'$ is at most the sum for $m$,
which is at most $B$, so $m'$ is in the window too.
\end{proof}

\begin{corollary}[Late siblings]\label{cor:late}
If every sibling message is newer than the ask node and the required fact sits
in message $m\in\A$, then $m\in C_\win$ if and only if
$\w(\Sib)\le\w(\{u\in\A:\tau(u)<\tau(m)\})$.
\end{corollary}
\begin{proof}
With late siblings, $\{u:\tau(u)\ge\tau(m)\}=\Sib\ \dot\cup\ \{u\in\A:\tau(u)\ge\tau(m)\}$.
Applying Proposition~\ref{prop:member} and $B=\w(\A)$:
$\w(\Sib)+\w(\A)-\w(\{u\in\A:\tau(u)<\tau(m)\})\le\w(\A)$, which is the stated condition.
\end{proof}

\begin{proposition}[Budget use]\label{prop:ratio}
If $C_\win\neq V$, then $1-\dfrac{w_{\max}}{B}<\dfrac{\w(C_\win)}{B}\le 1$.
\end{proposition}
\begin{proof}
The upper bound is the budget constraint. By maximality,
$\w(C_\win)>B-\w(v_{m+1})\ge B-w_{\max}$; divide by $B$.
\end{proof}

\section{Predictions checked against the data}

\begin{center}
\begin{tabular}{@{}lll@{}}
\toprule
Statement & Prediction & Observed (full run) \\
\midrule
Prop.~\ref{prop:eff} & $E=\w(\Sib)/\w(V)$ & exact in all 720 scenarios \\
Prop.~\ref{prop:nosib} & no canaries under ancestor-only & 0 of 2{,}598 answers \\
Prop.~\ref{prop:member} & which facts the window keeps & correct in 720 of 720 \\
Cor.~\ref{cor:root} & root fact never in window & 0 of 360 root cells kept it \\
Cor.~\ref{cor:late} & late siblings: kept iff $\w(\Sib)\le$ older path & correct in 360 of 360 \\
Prop.~\ref{prop:ratio} & window/budget ratio $>1-w_{\max}/B$ & mean 0.94, min 0.62 \\
\bottomrule
\end{tabular}
\end{center}

\end{document}
